\input{../../../assets/programming-editorial.tex}
\renewcommand{\studiumcoverphoto}{../../../assets/covers/vol09.jpg}
\urldef{\studiumsourceurl}\url{https://commons.wikimedia.org/wiki/File:IBM_1800_computer_at_Exxon_Research_and_Engineering_Company_laboratory_1971.jpg}
\urldef{\studiumlicenseurl}\url{https://commons.wikimedia.org/wiki/Commons:Public_domain}
\newcommand{\studiumcustomtitle}{
\studiumcover{Ingeniería de software}{Pruebas, arquitectura y proyectos reproducibles}{VOLUMEN 9}
\chapter*{Créditos y alcance de esta revisión}
Daniel Miguel Tejedor. Biblioteca de Programación, Studium. 2026.

Esta copia conserva los 75 capítulos de la colección original. La revisión de fuentes, cobertura y soluciones sigue abierta. Un cálculo aritmético reproducido no verifica un programa. Esta edición no cuenta con aprobación académica independiente.

\noindent\textbf{Fotografía de portada.} Courtesy of Science History Institute,  No Known Copyright / Wikimedia Commons / Public domain. IBM 1800 computer at Exxon Research and Engineering Company laboratory 1971.jpg.\par\studiumsourceurl

\noindent\textbf{Licencia y transformaciones.} Public domain.\par\studiumlicenseurl

Wikimedia thumbnail resized to 1920 pixels wide; derivative bytes preserved. Proportional placement and crop in shared LaTeX cover. La imagen y su recorte conservan la licencia indicada.

\noindent\textbf{Cómo estudiar.} Escribe los ejemplos en archivos, ejecuta las pruebas indicadas y compara la salida observada con la esperada. Consulta las soluciones después de intentarlo. Los bloques Terminal contienen comandos; Output muestra resultados y Diagnostic contiene errores.\par\clearpage
}
